\begin{solucion}
    % Añada este bloque justo después del resultado anterior.
% No introduce subsecciones ni numeración nueva: solo texto y ecuaciones.

\bigskip
Señalemos ahora un segundo hecho clásico:

\[
\sum_{\ell=0}^{n}\frac{D_\ell}{\ell!\,(n-\ell)!}=1,
\]

donde $D_\ell$ denota, como antes, el número de desarreglos de $\ell$ elementos.  
Proporcionamos una demostración puramente combinatoria contando de dos maneras el
conjunto $S_n$ de todas las permutaciones de $[n]$.

\textbf{(1) Conteo directo.}  
Sabemos que $\vert S_n\vert=n!$.

\textbf{(2) Clasificación según el número de puntos fijos.}  
Sea $\sigma\in S_n$ y denote‐\!$se$ por $\ell$ la cantidad de índices que
$\sigma$ deja fijos.  
Para reconstruir $\sigma$ procedemos en dos etapas independientes:

\begin{itemize}
  \item Elegimos los $\ell$ puntos fijos: existen $\binom{n}{\ell}$ maneras.
  \item Derangamos los $n-\ell$ elementos restantes, esto es, los
        permutamos sin puntos fijos; hay $D_{\,n-\ell}$ posibilidades.
\end{itemize}

En consecuencia, dado un $\ell$ determinado, el número de permutaciones con
exactamente $\ell$ puntos fijos es  
$\binom{n}{\ell}\,D_{\,n-\ell}$.

\medskip
Como toda permutación tiene algún número $\ell$ de puntos fijos comprendido entre
$0$ y $n$, sumando sobre todas las clases obtenemos

\[
n!\;=\;\sum_{\ell=0}^{n}\binom{n}{\ell}\,D_{\,n-\ell}.
\]

Basta ahora reindexar la suma (sustituyendo $\ell\mapsto n-\ell$) para
obtener la identidad equivalente

\[
n!\;=\;\sum_{\ell=0}^{n}\binom{n}{\ell}\,D_\ell.
\]

Finalmente dividimos ambos miembros por $n!$ y utilizamos la igualdad
$\displaystyle\binom{n}{\ell}=\frac{n!}{\ell!\,(n-\ell)!}$ para concluir

\[
\sum_{\ell=0}^{n}\frac{D_\ell}{\ell!\,(n-\ell)!}=1,
\]

tal como se quería demostrar.
\end{solucion}